\subsection{余均衡子的迷向群}

设 G 与 H 为群胚；设 \(\varphi\) 与 \(\psi\) 为从 H 到 G 的群胚态射。本节的目的在于总结上同伦子的迷向群的计算（II, 第 193 页, 命题 6）以及上一节中所作的上同伦子与余均衡子的迷向群的比较，以便由此推出余均衡子 Coeg( \(\varphi\) , \(\psi\) ) 的迷向群的计算。

以 \(\alpha\colon\mathrm{G}\to\mathrm{Coh}(\varphi,\psi)\) 、\(\beta\colon\mathrm{Coh}(\varphi,\psi)\to\mathrm{Coeg}(\varphi,\psi)\) 与 \(\gamma=\beta\circ\alpha\colon\mathrm{G}\to\mathrm{Coeg}(\varphi,\psi)\) 表示典范群胚态射。以 \(h\colon\mathrm{Som}(\mathrm{H})\to\mathrm{Fl}(\mathrm{Coh}(\varphi,\psi))\) 表示典范同伦；回忆 \(\mathrm{Coh}(\varphi,\psi)\) 的顶点集等于 \(\mathrm{Som}(\mathrm{G})\) 。

以 \(\Gamma\) 表示箭图 (Som(G), Som(H), Som(φ), Som(ψ))，再以 \(\varphi_0\) 与 \(\psi_0\) 表示由 \(\varphi\) 与 \(\psi\) 通过取轨道诱导的映射，以 \(\Gamma_0\) 表示 (Orb(G), Orb(H), \(\varphi_0, \psi_0\) ) ，即偶 ( \(\varphi, \psi\) ) 的框架。我们将假设箭图 \(\Gamma_0\) 是连通的且其顶点集非空，这等价于假设群胚 Coh(φ, ψ) 与 Coeg(φ, ψ) 都是传递的。

回忆（参看 II, 第 192 页, 定义 4），基本配备由下列数据组成：一个族 \((a, b, c_1, c_2, \mathrm{T}, i_0)\) ，其中：对所有 \(i \in \mathrm{Orb}(\mathrm{G})\) ，\(a(i)\) 是轨道 \(i\) （\(\mathrm{G}\) 的）中的一个顶点；对每个 \(j \in \mathrm{Orb}(\mathrm{H})\) ，\(b(j)\) 是轨道 \(j\) （\(\mathrm{H}\) 的）中的一个顶点，\(c_1(j)\) 与 \(c_2(j)\) 分别是 \(\mathrm{G}\) 中从 \(\varphi(b(j))\) 到 \(a(\varphi_0(j))\) 以及从 \(\psi(b(j))\) 到 \(a(\psi_0(j))\) 的箭；\(\mathrm{T}\) 是 \(\Gamma_0\) 的一个子箭图，其相伴树是图 \(\widetilde{\Gamma}_0\) 的一个极大树；最后，\(i_0\) 是 \(\mathrm{G}\) 的一个轨道，且我们令 \(a_0 = a(i_0)\) 。

我们随后定义一个箭图态射 \(\tau_0\) （从 \(\Gamma_0\) 到 \(\mathrm{Coh}(\varphi, \psi)\) ），使得 \(\tau_0(i) = \beta(a(i))\) 以及 \(\tau_0(j) = \alpha(c_1(j))^{-1}h(b(j))\alpha(c_2(j))\) （对 \(i \in \mathrm{Orb}(\mathrm{G})\) 及 \(j \in \mathrm{Orb}(\mathrm{H})\) ）。若 \(i\) 属于 \(\mathrm{Orb}(\mathrm{G})\) ，设 \(d_i\) 为图 \(\widetilde{\mathrm{T}}\) 中连接 \(i_0\) 与 \(i\) 的唯一的道路类，并以 \(\delta_i\) 表示它在 \(\mathrm{Coh}(\varphi, \psi)\) 中的、由群胚态射 \(\tau_0\) 给出的像。

由这些数据，II, 第 193 页的命题 6 给出一个满射同态。

\[
\Lambda \colon \bigl (\underset {i \in \pi_ {0} (\mathrm{G})} {*} \mathrm{G} _ {a (i)} \bigr) * \mathrm{F} (\operatorname{Orb} (\mathrm{H})) \to \operatorname{Coh} (\varphi , \psi) _ {a (i _ {0})}
\]

并描述其核的生成元，从而给出群 \(\mathrm{Coh}(\varphi,\psi)_{a_{0}}\) 的一个表现。

长度 \(\geqslant 1\) 的回路所成的集——在图 \(\widetilde{\Gamma}\) （与 \(\Gamma\) 相伴的）中——与集 \(\Omega(\widetilde{\Gamma})\) ——即序列 \((z_1, \ldots, z_n)\) （\(\widetilde{\Gamma}\) 的箭组成的，以 \(\mathbf{Z}/n\mathbf{Z}\) 为指标集）所成的集——等同，其中 \(n\) 取遍整数集 \(\geqslant 1\) ，且对每个 \(k \in \mathbf{Z}/n\mathbf{Z}\) ，\(z_k\) 的靶是 \(z_{k+1}\) 的起点。设 \(\mathbf{Z}\) 为 \(\Omega(\widetilde{\Gamma})\) 的一个子集，使共轭类 \(c(z)\) （对 \(z \in \mathbf{Z}\) ）生成满射同态 \(\beta_{a_0}\) 的核。同态 \(\beta_{a_0} \circ \Lambda\) 是满射的；为由此推出它的核，即群 \(\operatorname{Coeg}(\varphi, \psi)_{\beta(a_0)}\) 的一个表现，还需对每个 \(z \in \mathbf{Z}\) 选取一个元素 \(\mathrm{C}(z)\) （在群 \(\left( \begin{array}{cc} * & \mathrm{G}_{a(i)} \\ i \in \mathrm{Orb}(\mathrm{G}) & \end{array} \right) * \mathrm{F}(\pi_0(\mathrm{H}))\) 中），使 \(\Lambda(\mathrm{C}(z))\) 属于共轭类 \(c(z)\) 。

于是设 \(z = (z_{1}, \ldots, z_{n})\) 为 \(\Omega(\tilde{\Gamma})\) 的一个元素；令 \(z_{k} = (y_{k}, \varepsilon_{k})\) ，其中 \(y_{k} \in \mathrm{Fl}(\Gamma) = \mathrm{Som}(\mathrm{H})\) 且 \(\varepsilon_{k} \in \{\pm 1\}\) 。根据定义，\(c(z)\) 是元素的共轭类

\[
g h (y _ {1}) ^ {\varepsilon_ {1}} \dots h (y _ {n}) ^ {\varepsilon_ {n}} g ^ {- 1}
\]

的、在群 \(\mathrm{Coh}(\varphi,\psi)_{a_{0}}\) 中的共轭类，其中 g 是 \(\mathrm{Coh}(\varphi,\psi)\) 中连接 \(a_{0}\) 与 \(z_{1}\) 的起点的任意一条箭。对 \(k \in Z/nZ\) ，设 \(j_{k}\) 为 \(y_{k}\) 在 H 中的轨道，并选取 H 的一条箭 \(f_{k}\) ，它连接 \(y_{k}\) 与顶点 \(b(j_{k})\)。根据同伦的定义，我们随后有关系式

\[
h (y _ {k}) (\alpha \circ \psi) (f _ {k}) = (\alpha \circ \varphi) (f _ {k}) h (b (j _ {k}))
\]

在群胚 \(\mathrm{Coh}(\varphi,\psi)\) 中成立。于是，利用箭图态射 \(\tau_{0}\) 的定义，我们有

\[
\begin{array}{l} h (y _ {k}) = (\alpha \circ \varphi) (f _ {k}) \cdot h (b (j _ {k})) \cdot (\alpha \circ \psi) (f _ {k}) ^ {- 1} \\ \quad = (\alpha \circ \varphi) (f _ {k}) \cdot (\alpha \circ c _ {1}) (j _ {k}) \cdot \tau_ {0} (j _ {k}) \cdot (\alpha \circ c _ {2}) (j _ {k}) ^ {- 1} \cdot (\alpha \circ \psi) (f _ {k}) ^ {- 1} \\ \quad = (\alpha \circ \varphi) (f _ {k}) \cdot (\alpha \circ c _ {1}) (j _ {k}) \cdot \delta_ {\varphi_ {0} (j _ {k})} ^ {- 1} \cdot \\ \quad \cdot \delta_ {\varphi_ {0} (j _ {k})} \cdot \tau_ {0} (j _ {k}) \cdot \delta_ {\psi_ {0} (j _ {k})} ^ {- 1} \cdot \\ \quad \cdot \delta_ {\psi_ {0} (j _ {k})} \cdot (\alpha \circ c _ {2}) (j _ {k}) ^ {- 1} \cdot (\alpha \circ \psi) (f _ {k}) ^ {- 1} \\ = u _ {k} \Lambda (j _ {k}) v _ {k}, \end{array}
\]

其中我们已令

\[
u _ {k} = \alpha (\varphi (f _ {k}) c _ {1} (j _ {k})) \cdot \delta_ {\varphi_ {0} (j _ {k})} ^ {- 1} \quad \text { and } \quad v _ {k} = \delta_ {\psi_ {0} (j _ {k})} \cdot \alpha (\psi (f _ {k}) c _ {2} (j _ {k})) ^ {- 1}.
\]

对每个元素 \(k \in \mathbf{Z}/n\mathbf{Z}\) ，定义箭 \(\widetilde{u}_{k}, \widetilde{v}_{k}\) （G 中的）如下

\[
h (y _ {k}) ^ {\varepsilon_ {k}} = \widetilde {u} _ {k} \Lambda (j _ {k}) ^ {\varepsilon_ {k}} \widetilde {v} _ {k},
\]

使得

\[
(\widetilde {u} _ {k}, \widetilde {v} _ {k}) = \left\{ \begin{array}{l l} (u _ {k}, v _ {k}) & \text { if } \varepsilon_ {k} = 1; \\ (v _ {k} ^ {- 1}, u _ {k} ^ {- 1}) & \text { if } \varepsilon_ {k} = - 1. \end{array} \right.
\]

以 \(x_{k}\) 表示箭 \(h(y_{k})^{\varepsilon_{k}}\) 的起点；其靶则为 \(x_{k+1}\) ；设 \(i_{k}\) 为 \(x_{k}\) 的轨道。于是我们用下列公式定义群胚 G 中的回路 \(\lambda_{k}(z)\) （以 \(a(i_{k})\) 为基点的）：

\[
\lambda_ {k} (z) = \left\{ \begin{array}{l l} c _ {2} (j _ {k - 1}) ^ {- 1} \psi (f _ {k - 1}) ^ {- 1} \varphi (f _ {k}) c _ {1} (j _ {k}) & \text {if} (\varepsilon_ {k - 1}, \varepsilon_ {k}) = (1, 1); \\ c _ {2} (j _ {k - 1}) ^ {- 1} \psi (f _ {k - 1}) ^ {- 1} \psi (f _ {k}) c _ {2} (j _ {k}) & \text {if} (\varepsilon_ {k - 1}, \varepsilon_ {k}) = (1, - 1); \\ c _ {1} (j _ {k - 1}) ^ {- 1} \varphi (f _ {k - 1}) ^ {- 1} \varphi (f _ {k}) c _ {1} (j _ {k}) & \text {if} (\varepsilon_ {k - 1}, \varepsilon_ {k}) = (- 1, 1); \\ c _ {1} (j _ {k - 1}) ^ {- 1} \varphi (f _ {k - 1}) ^ {- 1} \psi (f _ {k}) c _ {2} (j _ {k}) & \text {if} (\varepsilon_ {k - 1}, \varepsilon_ {k}) = (- 1, - 1) \end{array} \right.\tag{1}
\]

根据构造，我们有

\[
\Lambda (\lambda_ {k} (z)) = \widetilde {v} _ {k - 1} \widetilde {u} _ {k},
\]

从而元素在同态 \(\Lambda\) 下的像

\[
\mathrm{C} (z) = \lambda_ {1} (z) (j _ {1}) ^ {\varepsilon_ {1}} \lambda_ {2} (z) (j _ {2}) ^ {\varepsilon_ {2}} \dots \lambda_ {n} (z) (j _ {n}) ^ {\varepsilon_ {n}}\tag{2}
\]

属于共轭类 \(c(z)\) 。

定义 3. —— 设 G 与 H 为群胚；设 \(\varphi\) 与 \(\psi\) 为从 H 到 G 的群胚态射。我们假设群胚 Coeg \((\varphi, \psi)\) 是传递的。设 \((a, b, c_1, c_2, T, i_0)\) 为偶 \((\varphi, \psi)\) 的一个基本配备。

我们称补配备为下列数据：一个子集 Z （\(\Omega(\widetilde{\Gamma})\) 的），使共轭类 \(c(\mathbf{z})\) （对 \(\mathbf{z} \in \mathbf{Z}\) ）生成同态 \(\beta_{a(i_0)}\) 的核；并且，对 Z 的每个元素 \(\mathbf{z} = ((y_1, \varepsilon_1), \ldots, (y_n, \varepsilon_n))\) ，给出箭的一个序列 \(f(\mathbf{z}) = (f_1, \ldots, f_n)\) （H 中的），使得 \(f_k\) 连接 \(y_k\) 与顶点 \(b(j_k)\) ，其中 \(j_k\) 表示 \(y_k\) 在 H 中的轨道。

偶 \((\varphi,\psi)\) 的一个完备配备是指一个基本配备与一个补配备的数据。

命题 5. —— 设 G 与 H 为群胚；设 \(\varphi\) 与 \(\psi\) 为从 H 到 G 的群胚态射。假设群胚 \(\operatorname{Coeg}(\varphi, \psi)\) 是传递的。以 \(\gamma\colon \mathrm{G} \to \mathrm{Coeg}(\varphi, \psi)\) 表示典范群胚态射。

我们给偶 \((\varphi, \psi)\) 配备一个完备配备

\[
(a, b, c _ {1}, c _ {2}, \mathrm{T}, i _ {0}, \mathrm{Z}, (f (\mathbf {z})) _ {\mathbf {z} \in \mathrm{Z}}).
\]

对 \(j \in \text{Orb(H)}\) ，设 \(\varphi_j \colon \text{H}_{b(j)} \to \text{G}_{a(\varphi_0(j))}\) 与 \(\psi_j \colon \text{H}_{b(j)} \to \text{G}_{a(\psi_0(j))}\) 分别为群同态 \(\text{Int}(c_1(j))^{-1} \circ \varphi_{b(j)}\) 与 \(\text{Int}(c_2(j))^{-1} \circ \psi_{b(j)}\) ，其中 \(\varphi_0\) 与 \(\psi_0\) 表示由 \(\varphi\) 与 \(\psi\) 通过取轨道诱导的典范映射。设 \(\tau\) 为框架 \(\Gamma_0\) （偶 \((\varphi, \psi)\) 的）到 \(\text{Coeg}(\varphi, \psi)\) 中的态射，由 \(\text{Som}(\tau)(i) = \gamma(a(i))\) （若 \(i \in \text{Orb(G)}\) ）定义，且使 \(\text{Fl}(\tau)(j)\) 是道路 \(\gamma(c_1(j))^{-1} \gamma(c_2(j))\) （在 \(\text{Coeg}(\varphi, \psi)\) 中的）。若 \(i\) 是 G 的一个轨道，设 \(c_i\) 为 \(\widetilde{\text{T}}\) 中连接 \(i_0\) 与 \(i\) 的唯一的道路类，并令 \(\delta_i = \widetilde{\tau}(c_i)\) ，其中 \(\widetilde{\tau} \colon \text{Grp(G)} \to \text{Coeg}(\varphi, \psi)\) 是由 \(\tau\) 诱导的典范群胚态射。

于是存在唯一的群同态

\[
\lambda \colon \bigl ( \begin{array}{c} * \\ i \in \operatorname{Orb} (\mathrm{G}) \end{array} \mathrm{G} _ {a (i)} \bigr) * \mathrm{F} (\operatorname{Orb} (\mathrm{H})) \to \operatorname{Coeg} (\varphi , \psi) _ {\gamma (a (i _ {0}))}
\]

使得

\[
\begin{array}{l l} \lambda (f) = \delta_ {i} \gamma_ {a (i)} (f) \delta_ {i} ^ {- 1} & \text { for } i \in \mathrm{Orb(G)} \text { and } f \in \mathrm{G} _ {a (i)}, \\ \lambda (j) = \delta_ {\varphi_ {0} (j)} \tau (j) \delta_ {\psi_ {0} (j)} ^ {- 1} & \text { for } j \in \mathrm{Orb(H)}. \end{array}
\]

同态 \(\lambda\) 是满射的；其核是包含下列元素的最小正规子群：

\[
\left(\mathrm{R} _ {1}\right) \quad r _ {1} (j) = j \quad \text { for } j \text { in } \mathrm{Fl} (\mathrm{T});
\]

\[
(\mathrm{R} _ {2}) \quad r _ {2} (j, f) = \varphi_ {j} (f) j \psi_ {j} (f) ^ {- 1} j ^ {- 1}
\]

\[
\text { for } j \in \operatorname{Orb} (\mathrm{H}) \text { and } f \in \mathrm{H} _ {b (j)};\tag{\((R_{3})\)}
\]

\[
r _ {3} (z) = \lambda_ {1} (z) j _ {1} ^ {\varepsilon_ {1}} \lambda_ {2} (z) j _ {2} ^ {\varepsilon_ {2}} \dots \lambda_ {n} (z) j _ {n} ^ {\varepsilon_ {n}}
\]

\[
\text { for } z = ((y _ {1}, \varepsilon_ {1}), \dots , (y _ {n}, \varepsilon_ {n})) \in \mathbf {Z},
\]

其中回路 \(\lambda_{i}(z)\) 由第 208 页的公式 (1) 定义。

这样一个态射的存在性与唯一性由群的族的自由积的泛性质得出（A, I, 第 85 页, 命题 8）。以 \(\alpha\colon G\to\mathrm{Coh}(\varphi,\psi)\) 与 \(\beta\colon\mathrm{Coh}(\varphi,\psi)\to\mathrm{Coeg}(\varphi,\psi)\) 表示典范态射，于是 \(\gamma=\beta\circ\alpha\)。再设 \(h\colon\mathrm{Som}(H)\to\mathrm{Fl}(\mathrm{Coh}(\varphi,\psi))\) 为联系 \(\alpha\circ\varphi\) 与 \(\alpha\circ\psi\) 的典范同伦。由于 \(\beta\) 是通过缩并 \(h\) 的像中的箭而得到的群胚态射，我们有

\[
\operatorname{Fl} (\tau) (j) = \gamma (c _ {1} (j)) ^ {- 1} \gamma (c _ {2} (j)) = \gamma (c _ {1} (j) ^ {- 1}) \beta (h (j)) \gamma (c _ {2} (j)) = \beta (\tau_ {0} (j))
\]

在 \(\operatorname{Coeg}(\varphi,\psi)\) 中成立，其中 \(\tau_{0}\colon\Gamma_{0}\to\mathrm{Coh}(\varphi,\psi)\) 表示由基本配备 \((a,b,c_{1},c_{2},\mathrm{T},i_{0})\) 诱导的箭图态射。于是，同态 \(\lambda\) 是命题 6（II, 第 193 页）中所定义的同态 \(\Lambda\) 与满射同态 \(\beta_{a(i_0)}\) 的复合。它特别是满射的。

设 \(z \in \mathbf{Z}\)。根据构造，\(\Lambda(r_3(z))\) 属于共轭类 \(c(z)\) 。根据补配备的定义，同态 \(\beta_{a(i_0)}\) 的核因此是 \(\mathrm{Coh}(\varphi, \psi)_{a(i_0)}\) 中包含元素 \(\Lambda(r_3(z))\) （对 \(z \in \mathbf{Z}\) ）的最小正规子群。由于同态 \(\Lambda\) 是满射的，同态 \(\lambda = \beta_{a(i_0)} \circ \Lambda\) 的核因此是群 \(\left( *_{i \in \mathrm{Orb}(G)} G_{a(i)} \right) * F(\pi_0(H))\) 中包含 \(\Lambda\) 的核的、由公式 \((R_1), (R_2)\) 给出的生成元以及由公式 \((R_3)\) 定义的元素的最小正规子群。命题得证。

